% 考古天文学（综述）
% license CCBYNCSA3
% type Wiki

本文根据 CC-BY-SA 协议转载翻译自维基百科\href{https://en.wikipedia.org/wiki/Archaeoastronomy}{相关文章}。

\begin{figure}[ht]
\centering
\includegraphics[width=6cm]{./figures/b9737bc319a7b6fe.png}
\caption{冬至之日，旭日东升，阳光仅在此时照亮爱尔兰纽格兰奇（Newgrange）墓冢的内部石室。} \label{fig_ArmY_1}
\end{figure}
考古天文学（Archaeoastronomy，又作 Archeoastronomy）是一门跨学科（interdisciplinary）[1]或多学科（multidisciplinary）[2]的研究领域，主要探讨古代人类如何“理解天空中的现象、如何利用这些现象，以及天空在其文化中所扮演的角色”[3]。Clive Ruggles 指出，将考古天文学简单地视为“古代天文学的研究”是误导性的，因为现代天文学是一门科学学科，而考古天文学则关注不同文化对天象的象征性与文化性解释[4][5]。考古天文学常与民族天文学（ethnoastronomy）并行研究，后者是一门人类学分支，专门研究当代社会中的观星活动与文化观念。此外，考古天文学还与历史天文学（historical astronomy）密切相关—— 后者利用古代的天象记录来解决天文学问题，以及天文学史（history of astronomy），即通过书面文献来评估古代天文学实践的发展历程[6]。
\begin{figure}[ht]
\centering
\includegraphics[width=6cm]{./figures/5d3ecc0dd8759825.png}
\caption{从西西里岛丰达凯利－凡蒂纳（Fondachelli Fantina）的史前遗址皮佐·文托（Pizzo Vento）观测到的春分或秋分日落景象。} \label{fig_ArmY_2}
\end{figure}
考古天文学（Archaeoastronomy）使用多种方法来揭示古代天文实践的证据，这些方法包括考古学、 人类学、 天文学、 统计学与概率论、以及历史学[7]。由于这些方法来源多样、数据类型差异巨大，如何将它们整合为一致而连贯的论证，一直是考古天文学家面临的长期挑战[8]。考古天文学在景观考古学（landscape archaeology）和认知考古学（cognitive archaeology）中占据互补地位。物质证据及其与天空的关联可以揭示更广阔的自然景观如何融入信仰体系，例如玛雅天文学（Mayan astronomy）与农业周期之间的关系[9]。其他结合认知与景观概念的研究实例包括：探讨定居点道路布局中所体现的宇宙秩序（cosmic order）[10][11]。  

考古天文学可适用于所有文化与各个历史时期。尽管不同文化对天空的意义理解差异显著，但在研究古代信仰时，仍然可以跨文化地应用一套科学方法[12]。或许正是因为必须在社会学维度与科学方法之间取得平衡，Clive Ruggles 才将考古天文学形容为—— “一个在一端产出高质量学术研究，而在另一端则存在近乎疯狂的无拘臆测的领域”[13]。
\subsection{历史}
在他所撰写的《天文考古学（Astro-archaeology）》简史中，约翰·米切尔（John Michell）指出，过去两个世纪以来，对古代天文学的研究地位经历了这样的转变：“从疯狂（lunacy）到异端（heresy），再到有趣的想法（interesting notion），最终到达正统的大门（the gates of orthodoxy）。”  

近二十年之后，我们仍可以追问：  
考古天文学是否依然停留在正统的大门前，  
还是已经迈入其间？  

—— 托德·博斯特威克（Todd Bostwick）引述约翰·米切尔[14]

两百年前，当约翰·米切尔尚未写下上述文字时，世界上尚无“考古天文学家”，也尚未出现“职业考古学家”，但已有天文学家与古物学家（antiquarians）。他们的一些研究工作被视为考古天文学的先驱：例如古物学家们注意到英格兰乡间遗迹的天文朝向，正如威廉·斯图克利（William Stukeley）在1740年对巨石阵（Stonehenge）所做的研究[15]；约翰·奥布里（John Aubrey，1678）[16]与亨利·尚西（Henry Chauncy，1700）[17] 则在教堂的朝向中寻找类似的天文学原则。十九世纪晚期，理查德·普罗克特（Richard Proctor）与查尔斯·皮亚齐·史密斯（Charles Piazzi Smyth）等天文学家研究了金字塔的天文朝向[18]。

“考古天文学”（Archaeoastronomy）这一术语由伊丽莎白·切斯利·贝蒂（Elizabeth Chesley Baity）于1973年提出（依循欧安·麦基〔Euan MacKie〕的建议）[19][20]。然而，作为一门研究主题，其历史或许更为久远，这取决于对“考古天文学”概念的定义。  

克莱夫·拉格尔斯（Clive Ruggles）[21]指出，十九世纪中期的海因里希·尼森（Heinrich Nissen）或许可以被视为最早的考古天文学家。罗尔夫·辛克莱尔（Rolf Sinclair）[22]认为，十九世纪末至二十世纪初的诺曼·洛克耶（Norman Lockyer）可被称为“考古天文学之父（father of archaeoastronomy）”。欧安·麦基（Euan MacKie）[23]则将其起源界定得更晚，他写道：“……考古天文学的诞生与现代繁荣，无疑应归功于亚历山大·汤姆（Alexander Thom）在1930年代至1970年代间于英国的研究工作。”
\begin{figure}[ht]
\centering
\includegraphics[width=6cm]{./figures/204a1de1a23caa92.png}
\caption{早期的考古天文学家在不列颠群岛进行巨石遗迹（Megalithic constructs）的测量研究，例如位于伦敦德里郡（County Londonderry）的奥格利什（Auglish）遗址，试图通过统计分析发现其中的规律性模式。} \label{fig_ArmY_3}
\end{figure}
在20世纪60年代，工程师亚历山大·汤姆（Alexander Thom）与天文学家杰拉尔德·霍金斯（Gerald Hawkins）的研究重新激发了人们对古代遗址天文学特征的兴趣。霍金斯曾提出巨石阵（Stonehenge）是“一台新石器时代的计算机”[24]。尽管霍金斯的主张大多被否定[25]，但汤姆的研究却不同凡响——他对不列颠群岛巨石遗址的测量结果，提出了当时广泛存在精确天文学实践的假说[26]。欧安·麦基（Euan MacKie）意识到汤姆的理论需要通过实地检验，于1970年与1971年在阿盖尔郡（Argyllshire）的金特劳（Kintraw）独立石遗址展开发掘，以验证汤姆预测的石碑上方山坡上是否存在观测平台。结果确实发现了一个人工平台——这一发现似乎印证了汤姆关于长线天文对齐（long alignment）的假说（金特劳被认为是一个精确的冬至观测点）。麦基随后又在艾莱岛（Islay）的卡尔图恩石圈（Cultoon stone circle）检验了汤姆的几何理论，同样得到了积极结果。麦基因此大体接受了汤姆的结论，并据此发表了关于不列颠的新史前学著作[27]。然而，克莱夫·拉格尔斯（Clive Ruggles）在重新评估汤姆的实地工作后认为，汤姆关于“高度精确天文学”的主张并未得到充分证据支持[28]。尽管如此，汤姆的影响依然深远。埃德温·C·克鲁普（Edwin C. Krupp）在1979年写道：“几乎凭一己之力，他确立了考古天文学实地研究与解释的标准，而他令人惊叹的成果在过去三十年间引发了广泛争议。”[29] 他的影响力至今不衰，对观测数据进行统计检验的做法仍然是考古天文学的重要研究方法之一[30][31]。
\begin{figure}[ht]
\centering
\includegraphics[width=6cm]{./figures/b822574214449a3d.png}
\caption{人们提出过这样的观点：玛雅遗址（如乌斯马尔 Uxmal）可能是按照天文对齐（astronomical alignments）的原则建造的。} \label{fig_ArmY_4}
\end{figure}
在“新大陆”（New World）的研究路径中，人类学家更充分地考虑了天文学在美洲原住民文明（Amerindian civilizations）中的作用，这一方法与旧大陆的研究方式存在显著差异。他们能够利用欧洲史前学所缺乏的资料来源，例如民族志记录[32][33]以及早期殖民者留下的历史文献。在安东尼·阿韦尼（Anthony Aveni）[34][35]的开创性研究启发下，新大陆的考古天文学家能够提出有关动机（motives）的解释，而在旧大陆，这类解释往往只能停留在推测层面。然而，由于新大陆研究侧重于历史资料，其关于天文观测“高精度”的某些论断，在与欧洲统计学导向的研究相比时显得相对薄弱[36]。

这一分歧在1981年国际天文学联合会（IAU）于牛津召开的会议上达到顶点[37]。会议参与者的方法论和研究问题差异极大，以至于会议论文被分为两卷分别出版[38][39]。尽管如此，会议被认为成功地促进了研究者之间的交流与合作，自那以后，牛津系列会议每隔四至五年便会在世界各地持续举办。随后的会议推动了考古天文学向更加跨学科的方向发展。研究者们开始尝试将考古学的情境性（contextuality）与天文学结合[40]——这大体描述了当代考古天文学的现状。如今的考古天文学不再仅仅致力于“证明古代天文学的存在”，而是进一步追问：人们为何会对夜空产生兴趣？
\subsection{与其他学科的关系}
“……考古天文学最令人喜爱的特质之一，  
就是它能让来自不同学科的学者彼此争论不休。” 
 
—— 克莱夫·拉格尔斯（Clive Ruggles）[41]

考古天文学长期以来被视为一门跨学科（interdisciplinary）的研究领域，它结合书面与非书面证据，以探究其他文化的天文学。因此，它可以被视为连接多种研究古代天文学路径的桥梁：包括天文考古学（astroarchaeology，现已废弃的术语，指通过古代建筑或地景的朝向推导天文信息的研究）、天文学史（history of astronomy，主要研究书面文献证据），以及民族天文学（ethnoastronomy，借助民族历史文献与当代民族志研究）[42][43]。

考古天文学作为一门跨学科领域的发展也体现在研究者的背景多样性上。该领域的研究由受训于不同学科的学者共同开展。近期多篇博士论文的作者将其研究描述为涉及考古学与文化人类学；也包括历史学的多个分支——如特定地区与时期的历史、科学史与宗教史；同时探讨天文学与艺术、文学、宗教之间的关系。只有极少数研究者将其工作直接归类为“天文学”，即便如此，也通常仅作为次要范畴[44]。

无论是实践中的考古天文学家，还是观察该领域的学者，他们对考古天文学的理解方式各不相同。部分研究者将其与科学史联系起来，关注某一文化如何观察自然，并建立概念体系以为这些观察赋予秩序[45]；另一些则关注推动特定历史人物采用某种天文学概念或技术的政治动机[46][47]。艺术史学者理查德·波斯（Richard Poss）采取了更为灵活的立场，认为北美西南地区的天文岩画（astronomical rock art）应当依据“西方艺术史与艺术批评的诠释学传统”进行解读[48]。然而，天文学家关注的问题有所不同 —— 他们希望为学生揭示本学科的可识别先驱，并尤为关心一个核心问题：如何确认某个遗址确实是出于天文目的而有意设计的[49]。

考古学界的反应（The reactions of professional archaeologists to archaeoastronomy）专业考古学界对考古天文学（archaeoastronomy）的反应始终褒贬不一。一些学者表现出难以理解甚至敌意的态度——这种敌意的形式多样，既包括主流考古学界将其视为“边缘学科”而拒绝接纳，也包括由于考古学家注重文化背景而早期考古天文学家偏重定量方法，导致双方沟通困难、理念分歧[50]。然而，近年来考古学界逐渐吸收了许多来自考古天文学的研究成果与洞见，并将其纳入考古学教材[51]。如前所述，一些学生甚至以考古天文学为主题撰写博士论文，显示出该领域在学术界的影响力正逐步扩大。

由于考古天文学家对这一学科的性质与边界意见分歧极大，他们甚至在“该如何命名”这一问题上存在争论。三大国际学术组织均将考古天文学视为研究文化与天文学关系的学科，采用“文化中的天文学”（Astronomy in Culture）或其译名作为正式术语。迈克尔·霍斯金（Michael Hoskin）认为，该领域的重要部分在于“事实收集”而非“理论建构”，因此建议将这一方向命名为“考古地形学”（Archaeotopography）[52]。克莱夫·拉格尔斯（Clive Ruggles）与尼古拉斯·桑德斯（Nicholas Saunders）则提出，应使用“文化天文学”（Cultural Astronomy）作为总称，以统一不同民俗天文学研究方法[53]。另一些学者认为，“天文学”一词并不准确，他们主张研究的核心是“宇宙观”（cosmology）而非“天文学”（astronomy），并建议采用西班牙语术语“cosmovisión”以避免“logos”的科学理性含义[54]。

当研究争论围绕方法论极端化时，人们常用一种“颜色代号”来区分不同学派 —— 这一命名方式源自第一次牛津会议（Oxford Conference）论文集的装帧颜色，该会议首次明确区分了不同的研究取向[55]。“绿色学派”（Green，旧大陆）考古天文学家主要依赖统计方法，  
常被批评忽视天文行为的文化与社会语境；而“棕色学派”（Brown，新大陆）学者则拥有丰富的民族志与历史资料，但在测量与统计分析方面常被指“过于随意”[56]。自20世纪90年代初以来，如何在这两种方法间寻求整合与平衡，一直是考古天文学界的重要讨论议题[57][58]。
\subsection{方法论}
“长期以来，我一直相信，这样的多样性必然需要一种包罗万象的理论来解释。  
但我现在认为，当初那样的想法实在太天真了。”

—— 斯坦尼斯瓦夫·伊瓦尼舍夫斯基（Stanislaw Iwaniszewski）[59]

考古天文学（archaeoastronomy）并没有唯一的研究方法。学者之间的分歧往往并非出现在“自然科学家”与“社会科学家”之间，而是取决于研究者所能获得的数据类型或研究地区的不同。在“旧大陆”（Old World），研究者往往仅能依靠遗址本身获取信息；而在“新大陆”（New World），则可结合民族志（ethnography）与历史记录来补充。这种地域性分化使得考古天文学在不同地区的发展路径独立而多样，其影响至今仍可在各地研究中观察到。总体而言，研究方法可大致分为两类，不过许多近年的项目已开始将两种方法结合使用，以实现更全面的解释框架。
\subsubsection{绿色考古天文学}
“绿色考古天文学”一名来源于《旧大陆考古天文学》（Archaeoastronomy in the Old World）一书的绿色封面[60]。这种方法主要以统计学为基础，特别适用于史前遗址研究—— 在这些遗址中，社会文化证据相较于有文字记载的历史时期较为稀少。该方法的基本技术由亚历山大·汤姆（Alexander Thom）在其对英国巨石遗址的广泛测量中建立。

汤姆的核心研究目标是检验史前人群是否掌握了高精度的天文学。他认为，通过“地平线天文学”（horizon astronomy），观察者能够以日为单位，精确估算一年中的特定日期。这种观测方法依赖寻找一个特定地点，在某一天太阳会刚好落入地平线的某个凹口或山间缺口中。常见的例子是：太阳被山峰遮挡，但在某个特定日期，太阳会在山的另一侧重新露出极小一部分，形成所谓的“双重日落”（double sunset）现象。下方的示意动画展示了一个假想遗址的两次日落情景：一次发生在夏至前一天，另一次则正好是夏至当天。在后者中，太阳出现了典型的“双重日落”现象。

为了验证这一设想，汤姆（Thom）对数百处石阵（stone rows）与石环（stone circles）进行了测量。单个排列的方向可能纯属偶然，但他计划通过统计显示，这些排列的整体分布并非随机，从而证明至少有一部分排列的朝向具有明确的天文意图。他的研究结果显示，古人或许将一年划分为八、十六，甚至三十二个大致相等的时间区段。其中包括两个至日（solstices）、两个分点（equinoxes），以及四个“交叉季度日”（cross-quarter days）—— 即位于至日与分点之间的中点日。这些节气日期与中世纪凯尔特历法（Celtic calendar）中的节日相对应。[61]虽然并非所有这些结论都被后世学者完全接受，但汤姆的研究对考古天文学（archaeoastronomy）产生了深远而持久的影响，尤其是在欧洲地区。

尤安·麦基（Euan MacKie）支持汤姆的分析，他通过将新石器时代的英国与玛雅文明进行比较，为汤姆的研究加入了考古学背景，提出当时的社会可能是分层结构的。[27]为了验证这些假设，麦基在苏格兰的几处史前观测台遗址进行了发掘。例如金特劳（Kintraw）遗址以一块四米高的竖石（standing stone）而闻名。汤姆认为，这块石头正好指向远处地平线上位于尤拉岛（Jura）上两座山 —— 贝恩·夏尼德（Beinn Shianaidh）与贝恩·欧·凯利亚斯（Beinn o’Chaolias）——之间的缺口。[62]他推测，在冬至时，太阳会在这一地平线缺口处出现“双重日落”（double sunset）现象。 然而，从地面高度观察，这一景象会被地形中的山脊遮挡，  
观测者必须抬高约两米的高度——因此需要另一个观测平台。  
麦基在峡谷对岸确实发现了一个由小石块堆成的平台。  
虽然缺乏文物遗迹令一些考古学家质疑其用途，  
而岩相分析（petrofabric analysis）也未能得出决定性结论，  
但随后在梅斯豪（Maes Howe）[63]与布什丘陵金菱形饰板（Bush Barrow Lozenge）[64]上的进一步研究  
让麦基得出结论：虽然用“科学”（science）一词来形容史前观测或许过于时代错置，  
但汤姆关于“高精度对齐”的整体观点基本正确。[65]

\bigskip
与此相对，克莱夫·拉格尔斯（Clive Ruggles）则认为汤姆的调查在数据选择上存在问题。[66][67]  
也有学者指出，“地平线天文学”的精度受到近地平线折射变化的限制。[68]  
对“绿色考古天文学”（Green archaeoastronomy）更深层的批评在于：  
虽然它能回答“古人是否对天文学感兴趣”，  
但由于缺乏社会文化层面的分析，  
它难以解释“古人为何对此感兴趣”。  
正如人类学家基思·金蒂格（Keith Kintigh）所言：  
“直白地说，在许多情况下，某一考古天文学主张是否正确，  
对人类学的发展并无太大意义，因为这些信息无法回应当下的解释性问题。”[69]  

尽管如此，“对齐研究”（study of alignments）仍然是考古天文学的重要组成部分，  
尤其在欧洲地区，依然是该领域研究的核心内容之一。[70]
```

